\documentclass[11pt, letterpaper]{article}

% ── Packages ──────────────────────────────────────────────────────────────────
\usepackage[margin=1.1in]{geometry}
\usepackage{amsmath, amssymb}
\usepackage{graphicx}
\usepackage{xcolor}
\usepackage{hyperref}
\usepackage{booktabs}
\usepackage{array}
\usepackage{caption}
\usepackage{subcaption}
\usepackage{enumitem}
\usepackage{listings}
\usepackage{titlesec}
\usepackage{parskip}
\usepackage{lmodern}
\usepackage{microtype}
\usepackage{float}
\usepackage{tikz}
\usepackage{pgfplots}
\pgfplotsset{compat=1.18}

% ── Colors (SemiAnalysis-ish palette) ─────────────────────────────────────────
\definecolor{accent}{HTML}{1A56DB}       % blue headings
\definecolor{lightgray}{HTML}{F5F5F5}
\definecolor{codebg}{HTML}{F0F4FF}
\definecolor{nvidia}{HTML}{76B900}
\definecolor{amd}{HTML}{ED1C24}
\definecolor{callout}{HTML}{EFF6FF}
\definecolor{calloutborder}{HTML}{3B82F6}

% ── Typography ────────────────────────────────────────────────────────────────
\hypersetup{
    colorlinks=true,
    linkcolor=accent,
    citecolor=accent,
    urlcolor=accent,
}
\titleformat{\section}{\large\bfseries\color{accent}}{}{0em}{}[\titlerule]
\titleformat{\subsection}{\normalsize\bfseries}{}{0em}{}
\titleformat{\subsubsection}{\normalsize\itshape}{}{0em}{}
\setlength{\parskip}{6pt}

% ── Code style ────────────────────────────────────────────────────────────────
\lstset{
    backgroundcolor=\color{codebg},
    basicstyle=\small\ttfamily,
    breaklines=true,
    frame=single,
    rulecolor=\color{accent!30},
    xleftmargin=8pt, xrightmargin=8pt,
}

% ── Callout box ───────────────────────────────────────────────────────────────
\newenvironment{keypoint}{%
  \begin{center}\begin{minipage}{0.93\textwidth}
  \setlength{\fboxsep}{10pt}
  \colorbox{callout}{\begin{minipage}{\dimexpr\textwidth-2\fboxsep\relax}
  \color{black}\small
}{%
  \end{minipage}}\end{minipage}\end{center}
}

% ── Figure helper ─────────────────────────────────────────────────────────────
% Usage: \fig{filename}{caption}{label}
\newcommand{\fig}[3]{%
  \begin{figure}[H]
    \centering
    \includegraphics[width=0.96\textwidth]{figures/#1}
    \caption{#2}
    \label{fig:#3}
  \end{figure}
}

% ── Title block ───────────────────────────────────────────────────────────────
\title{%
  \vspace{-1.5cm}
  {\color{accent}\Large\textbf{MoE Roofline Oracle:}}\\[4pt]
  {\large A First-Principles Analytical Framework for\\
  Mixture-of-Experts Inference Hardware Selection}
}
\author{Kai Maeda}
\date{March 2026}

% ══════════════════════════════════════════════════════════════════════════════
\begin{document}
\maketitle
\vspace{-0.5cm}
\hrule height 1.5pt
\vspace{0.4cm}

% ── Abstract ──────────────────────────────────────────────────────────────────
\begin{abstract}
\noindent
Selecting the right GPU for large-scale Mixture-of-Experts (MoE) inference is
non-trivial. The same silicon can exhibit wildly different cost-per-token
depending on model architecture, parallelism strategy, sequence length, and
operator mix. Existing hardware comparison tools treat these dimensions
independently or rely on empirical benchmarks that are often unavailable for
pre-release hardware. We present the \textbf{MoE Roofline Oracle}: an
interactive, first-principles analytical simulator that applies the roofline
performance model \cite{williams2009roofline} jointly across prefill and decode
phases, parallelism configurations (TP, PP, EP), and a full TCO model to
produce Pareto-optimal hardware recommendations. The tool covers nine GPU
families from NVIDIA Hopper through Blackwell Ultra and AMD CDNA3/CDNA4, and
includes speculative decoding variants. All results are theoretical upper
bounds; empirical validation is left to deployment.
\end{abstract}

\tableofcontents
\newpage

% ══════════════════════════════════════════════════════════════════════════════
\section{Introduction}
% ══════════════════════════════════════════════════════════════════════════════

The frontier of large language model deployment has shifted decisively toward
Mixture-of-Experts architectures. DeepSeek-V3 (671B total, 37B active),
Mixtral-8x22B, and Llama-4's MoE variants have demonstrated that conditional
computation---routing each token through only a sparse subset of
experts---enables models to scale total parameter count while keeping per-token
FLOPs tractable. However, this sparsity fundamentally changes the
hardware bottleneck profile relative to dense transformers.

\begin{keypoint}
\textbf{The core tension in MoE inference:} Expert parallelism (EP) reduces
memory-per-GPU but introduces AlltoAll communication latency. Tensor
parallelism (TP) accelerates GEMM on large batches but adds AllReduce overhead
per layer. Decode is almost always memory-bandwidth-bound regardless of GPU
generation; prefill is compute-bound at scale. Getting hardware selection
wrong means paying for compute you cannot use or bandwidth you cannot saturate.
\end{keypoint}

Prior art in this space---InferenceX, LLM-Perf, and Artificial Analysis---
provides empirical benchmarks on a subset of GPU/model combinations. These
are invaluable, but they cannot answer the question ``what will a B300 SXM
cost to operate with DeepSeek-V3 at 200ms ITL SLA?'' before the hardware is
available for benchmarking.

The MoE Roofline Oracle addresses this gap by deriving performance bounds
analytically from first principles, with hardware parameters sourced from
vendor white papers and community benchmarks.

\fig{overview_dashboard}{MoE Roofline Oracle dashboard. Left: per-operator arithmetic intensity and roofline chart. Right: Pareto frontier of hardware configurations across ITL and cost-per-million-token axes.}{overview}

% ══════════════════════════════════════════════════════════════════════════════
\section{Background: The Roofline Model for LLM Inference}
% ══════════════════════════════════════════════════════════════════════════════

The roofline model \cite{williams2009roofline} bounds achievable performance
as:

\begin{equation}
  P_\text{achieved} = \min\!\left(\pi,\; \beta \cdot I\right)
  \label{eq:roofline}
\end{equation}

\noindent where $\pi$ is peak compute throughput (FLOP/s), $\beta$ is peak
memory bandwidth (bytes/s), and $I$ is the \emph{arithmetic intensity}
(FLOPs per byte of HBM traffic). The ridge point $I^* = \pi / \beta$ (measured
in FLOPs/byte) separates two regimes:

\begin{itemize}[noitemsep]
  \item $I < I^*$: \textbf{memory-bandwidth-bound} --- performance scales with
    $\beta$; compute resources are idle.
  \item $I \geq I^*$: \textbf{compute-bound} --- performance scales with
    $\pi$; memory bandwidth is not the bottleneck.
\end{itemize}

\subsection{Why Prefill and Decode Live in Different Regimes}

The key insight for LLM inference is that prefill and decode have
\emph{fundamentally different} arithmetic intensities:

\textbf{Prefill} (processing the input prompt in parallel): The GEMM
computes $(\mathbf{X} \in \mathbb{R}^{B \times S \times d}) \times
(\mathbf{W} \in \mathbb{R}^{d \times d})$. With batch $B$ and sequence length
$S$:

\begin{equation}
  I_\text{prefill} \approx \frac{2 B S d^2}{B S d \cdot \texttt{dtype} + d^2 \cdot \texttt{dtype}}
    \;\approx\; \frac{B S}{B S + d / (2 d)} \cdot \frac{d}{\texttt{dtype\_bytes}}
\end{equation}

For large $S$, $I_\text{prefill} \to d / \texttt{dtype\_bytes}$ — typically
hundreds of FLOPs/byte on H100. This is firmly compute-bound.

\textbf{Decode} (generating one new token): Batch size for a single token is
$B=1$. The weight matrix $\mathbf{W}$ is streamed cold from HBM on every
forward pass:

\begin{equation}
  I_\text{decode} \approx \frac{2 d^2}{d^2 \cdot \texttt{dtype\_bytes}} = \frac{2}{\texttt{dtype\_bytes}}
\end{equation}

For BF16, $I_\text{decode} = 1$ FLOPs/byte --- far below any modern GPU's
ridge point (typically 100--300 FLOPs/byte). \textbf{Decode is almost always
memory-bandwidth-bound.}

\fig{roofline_chart}{Roofline chart for DeepSeek-V3 on H100 SXM. The prefill operator (large circle) sits firmly in the compute-bound regime; decode (small circle near origin) is deep in the memory-bound regime. The ridge point marks the boundary.}{roofline}

\subsection{MoE Specifics}

For MoE layers with $E$ total experts and top-$k$ routing, only $k$ experts
fire per token. This changes the arithmetic intensity:

\begin{align}
  \text{FLOPs}_\text{MoE} &= 2 k \cdot B S \cdot d_\text{model} \cdot d_\text{ffn}
    && \text{\cite{rajbhandari2022deepspeed}} \\
  \text{Bytes}_\text{MoE} &= k \cdot 2 d_\text{model} \cdot d_\text{ffn} \cdot \texttt{dtype\_bytes}
    && \text{(weight streaming, decode)}
\end{align}

The net intensity is still low during decode ($\sim 1$--$2$ FLOPs/byte),
since each token loads a fresh set of $k$ expert weight matrices.
During prefill, large batches can amortize weight loading across many tokens,
recovering a high-intensity regime.

% ══════════════════════════════════════════════════════════════════════════════
\section{System Architecture}
% ══════════════════════════════════════════════════════════════════════════════

The oracle is implemented as a Python library (\texttt{oracle/}) with a
Streamlit frontend (\texttt{app.py}). Figure~\ref{fig:architecture} shows the
data flow.

\fig{architecture_diagram}{Oracle system architecture. Hardware specs flow into the roofline engine; model FLOPs/bytes from the workload model are combined with parallelism overhead and TCO assumptions to produce Pareto frontiers.}{architecture}

\subsection{Workload Model (\texttt{oracle/workload.py})}

For each model architecture (dense or MoE), we compute per-operator FLOPs and
HBM bytes for both prefill and decode phases. The derivations follow
\cite{kaplan2020scaling, hoffmann2022training, dao2022flashattention,
narayanan2021efficient}.

\textbf{Attention FLOPs and bytes} (per layer, batch $B$, sequence $S$):
\begin{align}
  \text{FLOPs}_\text{attn} &= 4 B S d^2 + 4 B H S^2 d_h + 5 B H S S_\text{kv}
    \label{eq:attn_flops} \\
  \text{Bytes}_\text{attn} &= 4 d^2 \cdot \texttt{dtype} + 2 B H S_\text{kv} d_h \cdot \texttt{kv\_dtype}
    \label{eq:attn_bytes}
\end{align}

where $H$ is the number of attention heads, $d_h = d/H$ is the head dimension,
and $S_\text{kv}$ is the full context length attended over. The KV cache term
dominates bytes during decode.

\subsection{Roofline Engine (\texttt{oracle/roofline.py})}

For each (model, GPU, parallelism config, batch size) tuple:

\begin{enumerate}[noitemsep]
  \item Compute total FLOPs and bytes via \texttt{workload.py}.
  \item Apply effective hardware peaks: $\pi_\text{eff} = \pi \cdot \eta_\text{compute}$,
    $\beta_\text{eff} = \beta \cdot \eta_\text{bw}$, where $\eta \in [0,1]$
    are per-GPU software efficiency parameters sourced from community benchmarks.
  \item Compute bottleneck time: $t = \max(F / \pi_\text{eff},\; B / \beta_\text{eff})$.
  \item Report TTFT (prefill bottleneck) and ITL (single-token decode bottleneck).
\end{enumerate}

Model FLOP Utilization (MFU) is defined as:

\begin{equation}
  \text{MFU} = \frac{P_\text{achieved}}{\pi_\text{nominal}}
    = \frac{\min(\pi_\text{eff},\; \beta_\text{eff} \cdot I)}{\pi_\text{nominal}}
\end{equation}

Prefill MFU is high ($\sim$60--90\%) because prefill is compute-bound.
Decode MFU is structurally low ($\sim$5--30\%) because the GPU is
memory-bandwidth-bound and most tensor cores are idle.

\subsection{Parallelism Overhead (\texttt{oracle/parallelism.py})}

Three parallelism strategies are modeled:

\textbf{Tensor Parallelism (TP):} Requires two AllReduce operations per
transformer layer (after attention output projection and after FFN
down-projection \cite{narayanan2021efficient}). Communication time:

\begin{equation}
  t_\text{AR} \approx \frac{2(N-1)}{N} \cdot \frac{M}{B_\text{NVLink}} + \alpha
  \label{eq:allreduce}
\end{equation}

where $M = B \cdot S \cdot d_\text{model} \cdot \texttt{dtype}$ is the
message size and $B_\text{NVLink}$ is unidirectional NVLink bandwidth.
Parallelism efficiency:

\begin{equation}
  \eta_\text{TP} = \frac{t_\text{compute}}{t_\text{compute} + t_\text{AR}}
\end{equation}

\textbf{Expert Parallelism (EP):} Two AlltoAll collectives per MoE layer
(token dispatch + expert output gather \cite{lepikhin2021gshard}). For $EP$
ranks with $k$ active experts per token, each rank sends $B k / EP$ token
vectors of size $d_\text{model} \cdot \texttt{dtype}$ bytes.

\textbf{Load Imbalance:} Random top-$k$ routing means some EP ranks receive
more tokens than the mean. We model the expected maximum load using the
balls-into-bins order statistic \cite{mitzenmacher2001power}:

\begin{align}
  \sigma &= \sqrt{\bar{L}(1 - k/EP)},\quad \bar{L} = B k / EP \notag \\
  L_\text{max} &\approx \bar{L} + \sigma \sqrt{2 \ln EP} \label{eq:loadimbal}
\end{align}

\textbf{Pipeline Parallelism (PP):} Pipeline bubble fraction following
\cite{huang2019gpipe}:

\begin{equation}
  f_\text{bubble} = \frac{PP - 1}{PP - 1 + m}
\end{equation}

where $m$ is the number of microbatches. For single-request decode, $m=1$
and pipeline overhead is significant.

\fig{parallelism_breakdown}{Parallelism efficiency vs. TP degree for DeepSeek-V3 decode on H100 SXM (NVLink 4). Efficiency drops sharply beyond TP=4 as AllReduce latency dominates.}{parallelism}

\subsection{TCO Model (\texttt{oracle/cost.py})}

Cost-per-million output tokens:

\begin{equation}
  C = \frac{\text{TCO}_\text{annual}}{\text{tok/s}_\text{cluster} \cdot T_\text{year} \cdot \eta_\text{avail}} \cdot 10^6
  \label{eq:cost}
\end{equation}

where:
\begin{align}
  \text{TCO}_\text{annual} &= \frac{\text{CapEx}}{D} + \text{OpEx}_\text{power} + \text{OpEx}_\text{overhead} \\
  \text{CapEx} &= N_\text{GPU} \cdot (C_\text{GPU} + C_\text{server}/8) \cdot 1.05 \\
  \text{OpEx}_\text{power} &= \text{TDP} \cdot N_\text{GPU} \cdot \text{PUE} \cdot C_\text{kWh} \cdot 8760 \\
  \text{OpEx}_\text{overhead} &= \phi_\text{opex} \cdot \text{CapEx}
\end{align}

$D$ is the depreciation period (years), $C_\text{kWh}$ is the power cost,
$\phi_\text{opex}$ is the non-power overhead fraction, and $\eta_\text{avail}
= 0.97$ accounts for cluster downtime \cite{aws2023sla}.

Cluster utilization $\eta_\text{util} \in [0,1]$ is treated as a separate
user-controlled parameter representing the fraction of time the cluster is
actively serving requests (demand-side, not hardware-intrinsic). The effective
cost scales as $C / \eta_\text{util}$.

% ══════════════════════════════════════════════════════════════════════════════
\section{Hardware Database}
% ══════════════════════════════════════════════════════════════════════════════

Table~\ref{tab:gpus} summarizes the GPU configurations modeled.
Software efficiency parameters $\eta_\text{compute}$ and $\eta_\text{bw}$
represent achievable fractions of peak throughput with production inference
stacks (TRT-LLM, vLLM), sourced from Artificial Analysis,
MLCommons, and vendor guides.

\begin{table}[H]
\centering
\caption{GPU hardware parameters modeled in the oracle.}
\label{tab:gpus}
\resizebox{\textwidth}{!}{%
\begin{tabular}{lrrrrrrrr}
\toprule
\textbf{GPU} & \textbf{BF16 TFLOP/s} & \textbf{HBM BW (TB/s)} & \textbf{HBM (GB)} & \textbf{TDP (W)} & \textbf{GPU \$k} & $\eta_\text{compute}$ & $\eta_\text{bw}$ & \textbf{Ridge Pt.} \\
\midrule
H100 SXM        & 989  & 3.35 & 80  & 700   & 30  & 0.92 & 0.90 & 295 \\
H200 SXM        & 989  & 4.80 & 141 & 700   & 40  & 0.90 & 0.90 & 206 \\
B200 SXM        & 2,200 & 8.00 & 192 & 1,000 & 70  & 0.78 & 0.82 & 275 \\
B300 SXM        & 2,500 & 9.00 & 288 & 1,200 & 90  & 0.75 & 0.80 & 278 \\
GB200 NVL72     & 2,200 & 8.00 & 192 & 1,000 & 42  & 0.82 & 0.85 & 275 \\
GB300 NVL72     & 2,500 & 9.00 & 288 & 1,200 & 56  & 0.78 & 0.82 & 278 \\
MI300X          & 1,307 & 5.30 & 192 & 750   & 15  & 0.70 & 0.78 & 247 \\
MI325X          & 1,307 & 6.00 & 256 & 750   & 20  & 0.70 & 0.78 & 218 \\
MI355X          & 2,560 & 8.00 & 288 & 1,000 & 35  & 0.65 & 0.75 & 320 \\
\bottomrule
\end{tabular}%
}
\end{table}

\noindent\textbf{A note on software efficiency:} NVIDIA Hopper GPUs achieve
higher $\eta_\text{compute}$ (0.90--0.92) because the CUDA/TRT-LLM stack has
been optimized for Hopper for over two years. Blackwell GPUs launch with lower
efficiency (0.75--0.82) as the TRT-LLM Blackwell kernel library is still
maturing. AMD CDNA3 and CDNA4 carry a persistent penalty (0.65--0.78) from
ROCm's relative immaturity in GEMM kernel coverage versus cuBLAS.

% ══════════════════════════════════════════════════════════════════════════════
\section{The Pareto Frontier: Latency vs. Cost}
% ══════════════════════════════════════════════════════════════════════════════

For each (GPU $\times$ precision $\times$ TP $\times$ EP $\times$ PP $\times$
batch size) configuration, the oracle computes ITL, TTFT, and $/M tokens. The
Pareto frontier is defined as the set of configurations where no alternative is
simultaneously cheaper \emph{and} faster.

\begin{keypoint}
\textbf{Reading the Pareto plot:} Points on the frontier represent the
``best achievable'' tradeoffs. Moving left (faster ITL) requires accepting
higher cost. Moving down (cheaper) requires accepting slower ITL. Any point
not on the frontier is strictly dominated.
\end{keypoint}

\fig{pareto_itl}{Pareto frontier: ITL (ms) vs. cost per million output tokens for DeepSeek-V3, BF16, across all modeled GPU families. Each curve represents one GPU; points along each curve vary TP, EP, and batch size.}{pareto_itl}

\fig{pareto_ttft}{Pareto frontier: TTFT (ms) vs. cost per million tokens for DeepSeek-V3. TTFT is compute-bound and scales primarily with GPU TFLOP/s, so Blackwell GPUs dominate the low-latency end.}{pareto_ttft}

\subsection{Key Observations}

\textbf{MI300X dominates on cost at relaxed latency.} With the lowest GPU cost
(\$15k vs. \$30k for H100) and 192 GB HBM enabling fewer GPUs for model
sharding, MI300X achieves the lowest $/M tokens for ITL $> 50$ ms. The ROCm
efficiency penalty is outweighed by the hardware cost advantage at high batch
sizes.

\textbf{H100 dominates the mid-range.} The combination of mature software
efficiency (0.92), competitive HBM3 bandwidth, and moderate GPU cost makes
H100 the Pareto frontier champion for ITL targets in the 10--50 ms range.

\textbf{B200 and GB200 NVL72 dominate at strict latency.} For sub-10ms ITL,
the 2$\times$ higher HBM bandwidth and NVLink domain of Blackwell systems
enables smaller batches without sacrificing throughput. The NVL72 rack's
72-GPU NVLink domain eliminates IB bottlenecks for large TP/EP configurations.

\textbf{SpecDec variants shift all curves left-and-down.} Adding speculative
decoding compresses ITL by $\sim$2$\times$ (NVIDIA Hopper) to $\sim$2.5$\times$
(NVIDIA Blackwell) at no additional GPU cost. AMD benefits less ($\sim$1.5$\times$)
due to ROCm's less mature spec decoding stack.

\fig{pareto_specdec}{Effect of speculative decoding on the Pareto frontier for H100 SXM and B200 SXM. The +SpecDec variants shift curves toward lower ITL and lower cost simultaneously (higher throughput amortizes fixed CapEx).}{pareto_specdec}

% ══════════════════════════════════════════════════════════════════════════════
\section{Arithmetic Intensity Deep Dive}
% ══════════════════════════════════════════════════════════════════════════════

\fig{intensity_breakdown}{Arithmetic intensity by operator and phase for Mixtral-8x22B and DeepSeek-V3. FFN attention bars show per-layer FLOPs/byte. Decode intensities cluster near 1--2 FLOPs/byte regardless of model size.}{intensity}

\subsection{How Batch Size Moves the Roofline Point}

Increasing decode batch size $B$ increases both FLOPs and bytes proportionally
(each additional sequence adds one more token's worth of compute and KV cache
reads). However, the weight matrices are shared across all sequences in a
batch, so weight-loading bytes do not scale with $B$ for large $d$:

\begin{equation}
  I_\text{decode}(B) \approx \frac{2 B d^2}{d^2 \cdot \texttt{dtype} + 2 B d \cdot \texttt{dtype}}
    = \frac{2B}{1 + 2B/d} \xrightarrow{B \ll d} \frac{2B}{\texttt{dtype}}
\end{equation}

At $B = 64$, $d = 7{,}168$ (DeepSeek-V3), $I \approx 64$ FLOPs/byte ---
approaching (but not yet reaching) the H100 ridge point of $\sim$295
FLOPs/byte. This explains why decode only becomes compute-bound at very
large batch sizes.

\fig{intensity_vs_batch}{Arithmetic intensity of DeepSeek-V3 decode vs. batch size, with GPU ridge points overlaid. GPU shading shows which GPUs are memory-bound vs. compute-bound at each batch size.}{intensity_batch}

\subsection{MoE Expert Routing and the EP Tradeoff}

Expert Parallelism shards experts across $EP$ GPU ranks. For a model with $E$
total experts and top-$k$ routing, each GPU holds $E/EP$ experts. With $EP$
ranks:

\begin{itemize}[noitemsep]
  \item \textbf{Memory benefit:} each GPU holds $1/EP$ of the expert weights.
    Enabling larger batch sizes per GPU.
  \item \textbf{Communication cost:} $2 \times EP$ AlltoAll calls per MoE
    layer (dispatch + gather), each moving $B k / EP \cdot d_\text{model}$
    tokens per rank over IB.
  \item \textbf{Load imbalance:} With random routing, some ranks receive up to
    $\sqrt{2 \ln EP}$ standard deviations above mean load (Eq.~\ref{eq:loadimbal}).
\end{itemize}

\fig{ep_sweep}{Expert Parallelism sweep for DeepSeek-V3 on H100 SXM: effective throughput (tok/s) vs. EP degree. The optimal EP varies with batch size and IB bandwidth.}{ep_sweep}

% ══════════════════════════════════════════════════════════════════════════════
\section{TCO Sensitivity Analysis}
% ══════════════════════════════════════════════════════════════════════════════

\subsection{Power Cost and PUE}

The cost-per-token is a roughly linear function of power cost $C_\text{kWh}$
for power-hungry GPUs (B200 at 1000W TDP vs. H100 at 700W). At
\$0.10/kWh (premium colocation), B200's power disadvantage closes the gap
with H100 by $\sim$\$0.10/M tokens.

\fig{sensitivity_power}{Cost sensitivity to power cost (\$/kWh) for H100 SXM, B200 SXM, and MI300X. MI300X benefits disproportionately from low-power regions due to lower TDP and higher HBM capacity per dollar.}{sensitivity_power}

\subsection{Cluster Utilization}

Cluster utilization $\eta_\text{util}$ represents the fraction of time the
cluster is actively serving requests. This is a demand-side input, not a
hardware property. Cost scales as $C \propto 1/\eta_\text{util}$: a cluster
running at 30\% utilization pays 3.3$\times$ more per token than one at
100\%.

\begin{keypoint}
\textbf{Practical implication:} For a startup serving intermittent traffic,
$\eta_\text{util} \approx 0.20$--$0.30$. For a high-traffic consumer product,
$\eta_\text{util} \approx 0.60$--$0.80$. The hardware that is cheapest at
100\% utilization may not be cheapest when your cluster sits idle 70\% of the
day (e.g., due to smaller per-GPU minimums for spot provisioning).
\end{keypoint}

\subsection{Depreciation and Refresh Cycle}

Blackwell hardware commands a faster refresh cycle (2--2.5 years estimated)
due to higher per-unit prices and rapid generational iteration. This
amortization difference shifts B200's effective CapEx contribution upward
relative to H100's standard 3-year cycle.

% ══════════════════════════════════════════════════════════════════════════════
\section{Model Catalog}
% ══════════════════════════════════════════════════════════════════════════════

Table~\ref{tab:models} lists the models included in the oracle.

\begin{table}[H]
\centering
\caption{Model configurations in the oracle catalog.}
\label{tab:models}
\begin{tabular}{lrrrrrrr}
\toprule
\textbf{Model} & \textbf{Total Params} & \textbf{Active Params} & $d$ & \textbf{Layers} & \textbf{Experts} & \textbf{Top-}$k$ \\
\midrule
Mixtral-8x7B     & 46.7B  & 12.9B & 4,096  & 32 & 8   & 2 \\
Mixtral-8x22B    & 141B   & 39.1B & 6,144  & 56 & 8   & 2 \\
DeepSeek-V2      & 236B   & 21B   & 5,120  & 60 & 160 & 6 \\
GPT-3 175B (dense) & 175B & 175B  & 12,288 & 96 & 1   & 1 \\
\bottomrule
\end{tabular}
\end{table}

DeepSeek-V2 is the most challenging case: 160 experts with top-6 routing
creates an unusually high EP communication demand, and the large expert count
means full expert weight loading per forward pass is extremely unlikely to fit
in any GPU's on-chip cache.

% ══════════════════════════════════════════════════════════════════════════════
\section{Limitations and Future Work}
% ══════════════════════════════════════════════════════════════════════════════

\subsection{Known Limitations}

\begin{enumerate}
  \item \textbf{Analytical upper bounds only.} The roofline model gives
    theoretical performance ceilings. Real deployments may be 10--30\% below
    these bounds due to kernel launch overhead, NCCL/RCCL library inefficiency,
    and OS scheduling jitter.

  \item \textbf{KV cache capacity not checked.} The oracle does not explicitly
    verify that the KV cache for a given batch size and context length fits in
    HBM. For long contexts ($> 32$k tokens) with large batches, this constraint
    can make some configurations infeasible.

  \item \textbf{Single speculative decoding speedup value.} The spec decode
    speedup is modeled as a fixed multiplier per GPU family. In reality,
    speedup depends on acceptance rate, draft model quality, and context.
    EAGLE-2 reports 2--3.5$\times$ depending on the task.

  \item \textbf{Homogeneous clusters only.} Prefill/decode disaggregation
    (e.g., Mooncake architecture with separate prefill and decode fleets on
    different hardware) is not yet modeled.

  \item \textbf{Pre-release hardware estimates.} B300 SXM and GB300 NVL72
    specs are estimated from announcements and architectural extrapolation.
    These are marked \texttt{specs\_confirmed=False} in the codebase.
\end{enumerate}

\subsection{Future Work}

\begin{itemize}[noitemsep]
  \item \textbf{Prefill/decode disaggregation:} Model separate fleet sizing for
    prefill-heavy vs. decode-heavy workloads (relevant for long-context chatbots
    and document QA).
  \item \textbf{KV cache offloading:} Model the regime where KV cache exceeds
    HBM capacity and spills to CPU DRAM or NVMe, with associated PCIe
    bandwidth penalty.
  \item \textbf{FP4 / MXFP4:} Blackwell and CDNA4 support native FP4 tensor
    cores. Modeling the quantization accuracy tradeoff and compute-efficiency
    gain is left to future work.
  \item \textbf{Empirical validation:} Cross-validate model predictions against
    MLCommons Inference results and Artificial Analysis benchmarks.
\end{itemize}

% ══════════════════════════════════════════════════════════════════════════════
\section{Conclusion}
% ══════════════════════════════════════════════════════════════════════════════

The MoE Roofline Oracle provides a principled, extensible framework for
hardware selection in large-scale MoE inference. By deriving performance
bounds from first principles and combining them with a full TCO model, it
enables apples-to-apples comparisons across GPU generations --- including
hardware not yet available for empirical benchmarking.

The key takeaways are:

\begin{enumerate}
  \item Decode is structurally memory-bandwidth-bound at small batch sizes.
    Choosing hardware with high HBM bandwidth-per-dollar (MI300X, H200)
    minimizes $/M tokens at relaxed latency targets.

  \item Prefill is compute-bound. For strict TTFT SLAs, TFLOP/s per dollar
    (Blackwell FP4) is the right optimization target.

  \item Parallelism efficiency drops sharply at high TP degrees, especially
    across nodes. The NVL72 rack-scale NVLink domain eliminates this penalty
    for intra-rack configurations.

  \item Speculative decoding disproportionately benefits compute-heavy GPUs
    (Blackwell), widening their advantage on decode throughput compared to
    memory-bandwidth-heavy GPUs (MI300X).

  \item Cluster utilization is the dominant uncertainty in cost modeling.
    Hardware selection should account for expected traffic patterns and
    idle time, not just peak throughput.
\end{enumerate}

% ══════════════════════════════════════════════════════════════════════════════
\bibliographystyle{plain}
\bibliography{refs}
% ══════════════════════════════════════════════════════════════════════════════

\end{document}
